\section{Herramientas Utilizadas}
El producto final del trabajo de investigación es el desarrollo de un add-on de Blender. 
Por ello, las herramientas a utilizar consisten principalmente de software para el
desarrollo.

\subsection{Visual Studio Code}
Visual Studio Code (o solo Code) es un editor de texto y ambiente de desarrollo integrado
(IDE) desarrollado por Microsoft. Code presenta varias características como soporte
para múltiples lenguajes de programación, herramientas de debug, integración con Git,
integración con la linea de comandos y extensiones que extienden la funcionalidad base 
del editor.
\medskip

Algunas de las extensiones que se usan incluyen Pylance, Python, Python Debugger y 
Python Environments (por Microsoft) para integrar funcionalidad específica de Python.

\subsection{Python}
Python es un lenguaje de programación orientado a objetos de uso general, siendo una
opción popular para desarrollo relacionado a inteligencia artificial. Además, es el
lenguaje de scripting usado por Blender.
\medskip

Python utiliza pip, un administrador de paquetería para desarrollo.
Permite la fácil instalación de bibliotecas y módulos como PyTorch y Numpy para su uso
en scripts de Python.
\medskip

Las librerías a usar incluyen:
\begin{description}
    \item[] \textit{Pytorch}, una biblioteca de código abierto desarrollado por Meta y sucesor
    de Torch. Pytorch provee funcionalidades y APIs para el desarrollo relacionado al
    aprendizaje profundo y constituye la base de muchas implementaciones de inteligencia
    artificial.
    \item[] \textit{Numpy}, una biblioteca matemática que implementa funcionalidad para
    arreglos matemáticos grandes y multidimensionales, así como su manipulación. Es
    una de las bibliotecas más utilizadas en el desarrollo científico con Python, y es usado
    en conjunto con Pytorch por su integración con los tensores.
    \item[] {Blender Python API (bpy)} es la bibliteca de Blender para Python que ofrece una
    API para interactuar con el motor gráfico y crear scripts para extender la funcionalidad
    de Blender.
    \item[] {pyopenvdb} es la implementación de OpenVDB para Python, permitiendo trabajar
    con los datos volumétricos en Python.
    \item[] {venv} es una utilidad de Python para utilizar ambientes virtuales. Los ambientes
    virtuales sirven de ambientes aislados del sistema con su propio interprete de Python
    y paquetes, funcionando de forma independiente a otras instalaciones y evitando
    problemas de compatibilidad relacionados a diferentes versiones.
\end{description}

\subsection{Blender}
Blender es una suite de modelado 3D de código abierto, desarrollado por la comunidad
y coordinado por la fundación Blender. Blender cuenta con un interprete de Python integrado
que permite la ejecución de scripts para extender su funcionalidad.
\medskip

Se utilizará la version 5.2 LTS de Blender por dos razones principales:
\begin{description}
    \item[] A partir de la versión 5.0 Blender integra en el editor de nodos las cuadrículas
    como un tipo de dato nativo, así como nuevos nodos para la manipulación de estas cuadrículas,
    lo cual en conjunto facilita en comparación a versiones anteriores flujos de trabajo
    procedural con datos volumétricos como OpenVDB.
    \item[] Al ser una versión LTS (Long Term Support, o soporte de largo plazo) se
    garantíza la estabilidad del programa así como su soporte a largo plazo.
\end{description}

Blender 5.2 LTS utiliza Python versión 3.13.13.

\subsection{Git}
Git es un software de control de versión altamente utilizado en el desarrollo de software.
Mediante el uso de ramas, permite a un desarrollador trabajar en distintas características de
un programa, de forma segura, sin que interfieran entre sí al realizar cambios. En consecuencia asegura
un desarrollo seguro de software al ser capaz de aislar diferentes versiones del código,
permitiendo separar versiones funcionales de las experimentales, así como revertir
cambios.
\medskip

GitHub es una plataforma en linea para alojar codigo, basado en Git, pero con un mayor
enfoque colaborativo y con herramientas extendidas para operaciones de desarrollo como
CI/CD. Git y GitHub suelen utilizarse en conjunto, con Git actuando como el administrador
de los cambios locales y Github actuando como repositorio remoto.

\subsection{Hardware}
Para el desarrollo del proyecto de investigación se utilizará una computadora personal con las
siguientes características:

\begin{description}
    \item[] Sistema operativo Windows 11 Home.
    \item[] Procesador Intel i7 de 12va generación, con velocidad base de 2.3 GHz.
    \item[] 16 GB de memoria RAM.
    \item[] Tarjeta gráfica NVIDIA RTX 3050 Ti (Laptop), con 4 GB de VRAM dedicada.
    \item[] 1TB de almacenamiento SSD
\end{description}

El uso de una tarjeta gráfica es particularmente beneficioso para ejecutar modelos de
inteligencia artificial, operaciones de convolución y programación grafica.